% Source-corrected proof, 2 October 2026; superseded v40 bytes preserved.

\begin{proof}[\Cref{thm:p031-orbit-fiber-maximum}]
Let $G$ be the group in
\Cref{def:p031-finite-analytic-orbit-presentation}, and put
\begin{equation}
\label{eq:p031-orbit-symmetrization}
                         m:=\max_{g\in G}u\circ g.
\end{equation}
Pullback by an analytic automorphism preserves piecewise linearity and
semipositivity, the latter by
\cite[Proposition~7.6\textup{(1')}]{GJR26}.

Suppose first that $u$ is semipositive $\mathbf Q$-PL.  Every translate
$u\circ g$ is semipositive $\mathbf Q$-PL, and $m$ is $\mathbf Q$-PL by
the finite-maximum property in \cite[\S5.2\textup{(7)}]{GJR26}.  The
translates are classically psh by
\cite[Theorem~4.3]{GR25Classical}; their finite maximum is classically psh
by \cite[Proposition~4.2\textup{(3)}]{GR25Classical}.  The converse
direction of that same Theorem~4.3 makes $m$ semipositive.

Suppose instead that $u$ is semipositive real PL and that $m$ is real PL.
The translates are classically psh by \cite[Theorem~4.3]{GR25Classical},
and their finite maximum is classically psh by
\cite[Proposition~4.2(3)]{GR25Classical}. The real-PL hypothesis permits
another application of Theorem~4.3, which makes $m$ semipositive.
No closure of the real-PL or LFS class under arbitrary maxima is used.

Equation \eqref{eq:p031-action-over-base} makes $m$ invariant under $G$,
and each fiber is one orbit.  There is therefore a unique function
$b:X\to\mathbf R$ satisfying
\begin{equation}
\label{eq:p031-orbit-symmetrization-descends}
                         m=b\circ\pi.
\end{equation}
In the rational case apply \cite[Proposition~5.11]{GJR26} with
$\Lambda=\mathbf Q$; in the conditional real-PL case apply it with
$\Lambda=\mathbf R$.  Since these coefficient groups are divisible and
$\pi$ is finite, flat, and surjective, the proposition shows that $b$ is
$\mathbf Q$-PL or real PL, respectively.  Semipositivity reflects through
the proper surjection $\pi$ by
\cite[Proposition~7.6\textup{(1')}]{GJR26}.

For $e\in\pi^{-1}(x)$, the orbit condition finally gives
\[
 b(x)=m(e)=\max_{g\in G}u(g(e))
              =\max_{e'\in\pi^{-1}(x)}u(e').
\]
Thus $b=\pi_{\max}u$.  In the second case the argument descends real
piecewise linearity and nef semipositivity, but constructs no semiample
residue model and hence proves no locally Fubini--Study conclusion.
\end{proof}

\begin{proof}[\Cref{thm:p031-orbit-envelope-comparison}]
Pullback preserves semipositive $\mathbf Q$-PL functions.  Hence every
base competitor gives a cover competitor and
\[
                  s_{\mathbf Q+,X}\circ\pi
                         \leq s_{\mathbf Q+,E}.
\]
Conversely, if $u\leq f\circ\pi$ is a semipositive $\mathbf Q$-PL cover
competitor, then \Cref{thm:p031-orbit-fiber-maximum} makes
$b:=\pi_{\max}u$ a semipositive $\mathbf Q$-PL base competitor.  Thus
\[
          u(e)\leq b(\pi(e))
              \leq s_{\mathbf Q+,X}(\pi(e)).
\]
Taking the supremum over the cover competitors proves the reverse
inequality in
\eqref{eq:p031-rational-semipositive-envelope-pullback}.  Pullback and
fiberwise maximization also show that the two rational families are empty
simultaneously.

Pullback of locally Fubini--Study base competitors gives
\[
                 s_{\mathrm{LFS},X}\circ\pi
                        \leq s_{\mathrm{LFS},E}.
\]
Every LFS cover competitor is a semipositive real-PL cover competitor,
so $s_{\mathrm{LFS},E}\leq s_{\mathrm{PL},E}$. Pullback of semipositive
real-PL base competitors gives
$s_{\mathrm{PL},X}\circ\pi\leq s_{\mathrm{PL},E}$. These prove
\eqref{eq:p031-lfs-pl-envelope-chain} without an additional maximum
hypothesis. If every cover LFS competitor has a real-PL orbit maximum,
the conditional second item of \Cref{thm:p031-orbit-fiber-maximum} makes
that maximum a semipositive real-PL base minorant dominating the original
competitor. Taking the supremum then gives the additional stated inequality.
\end{proof}

\begin{proof}[\Cref{prop:p031-real-pl-fiber-maximum-failure}]
Work over a complete nontrivially valued field whose additive value group
is $\mathbf Q$.  On $X=\mathbf G_{m,K}^{\mathrm{an}}$, Example~5.3 of
\cite{GJR26} gives the real-PL function
\[
                         h=\pi\log|t|+1
\]
for which $\max\{h,0\}$ is not real PL: under tropicalization it has its
unique bend at the irrational point $1/\pi$.  The function $h$ is LFS: choose $a\in K^\times$ with $-\log|a|=1$
and write
\[
 h=\pi(-\log|t^{-1}|)+(-\log|a|).
\]
This is the nonnegative real unit-logarithm presentation of
\cite[\S7.9]{GJR26}. In particular, $h$ is semipositive real PL.
The coefficient $\pi$ in this formula denotes the usual real constant.

Let $E=X\amalg X$, let $\pi_E:E\to X$ be the fold map, and let the group
of order two exchange the components.  This is a finite analytic orbit
presentation.  Define $u$ to equal $h$ on the first component and $0$ on
the second.  Then $u$ is LFS, hence semipositive real PL, whereas
\[
                         (\pi_E)_{\max}u=\max\{h,0\}
\]
is not real PL.  This proves the asserted failure.
\end{proof}
